El punto de partida del modelo es el conjunto de datos enriquecido obtenido en la sección \ref{SEC:PROC_DATASET}. A partir de este, se realiza un procesamiento final para definir las variables de entrada y salida del modelo.

En primer lugar, se define la característica \texttt{Service Port}, la cual se define como $\min(SP, DP)$, donde $SP$ es el puerto origen y $DP$ el puerto destino. El objetivo de esto fue reducir el ruido e introducir unicamente para el entrenamiento y evaluación el puerto del servicio, el cual suele corresponder al menor de los dos puertos involucrados en la comunicación.

A continuación, se establece la característica \texttt{Label} como variable objetivo o target. Esta contiene el tipo de tráfico al que pertenece el flujo, es decir, si es tráfico benigno o si pertenece a alguna clase de ataque. También se descartaron una serie de características para el entrenamiento: \texttt{Flow ID}, \texttt{Source IP}, \texttt{Destination IP} y \texttt{Timestamp} debido a que su función es de identificación; \texttt{Source Port} y \texttt{Destination Port} debido a que únicamente se va a utilizar el puerto del servicio; \texttt{Init\_Win\_bytes\_backward} e \texttt{Init\_Win\_bytes\_forward} debido a que identifican el Sistema Operativo del atacante y la víctima y pueden ser manipulados con facilidad por el atacante.

Finalmente, el conjunto de datos se dividió en un subconjunto de entrenamiento y otro de evaluación mediante la función \texttt{train\_test\_split} de la librería \texttt{scikit-learn}. Se decidió dividir la proporción de datos en los subconjuntos del 75\% para el subconjunto de entrenamiento y 25\% para el subconjunto de evaluación. El subconjunto de entrenamiento se utilizará para ajustar el modelo, mientras que el subconjunto de evaluación se utilizará para evaluar las predicciones del modelo. Para que no hubiera desbalance en las clases con pocas muestras, la división de estos subconjuntos se realizó de forma estratificada mediante el parámetro \texttt{stratify = y}. Además, se fijó \texttt{random\_state = 42} para que los resultados fueran los mismos entre distintas ejecuciones, ya que este parámetro fija la semilla del generador de números aleatorios.

\PythonCode{Separación en entrenamiento y evaluación}
{Fragmento de código utilizado para dividir el conjunto de datos en subconjuntos de entrenamiento y evaluación mediante \texttt{train\_test\_split}.}
{python/train_split.py}{1}{999}{1}